\chapter{What Is Programming Really?}

\section{Introduction: The Messy Desk Problem}

It is 4:00 PM on a Tuesday. Alex walks into the bedroom and drops a heavy backpack on the floor. On the desk, there is a messy pile of math homework, a physics assignment due tomorrow, and a sticky note reminding Alex to go to the grocery store for milk. At the same time, Maya sends a text message: \textit{"Hey Alex, want to play an online game?"}

Alex looks at the mess, feels overwhelmed, and thinks, \textit{"If I just start playing, my math homework won't get done. If I study first, I might miss out. I need a clear plan!"}

Without knowing it, Alex is about to become a programmer.

Programming is often thought of as writing secret code on a dark computer screen. But at its heart, programming is simply **problem-solving**. It is the art of taking a messy, confusing problem, breaking it down into tiny, orderly steps, and writing them down so clearly that even a machine can follow them without getting confused.

\begin{definition}[Programming]
Programming is the process of giving step-by-step instructions to a computer to solve a specific problem.
\end{definition}

\section{Planning the Afternoon: What is an Algorithm?}

To solve the afternoon chaos, Alex takes out a notebook and writes down a plan:
\begin{enumerate}
    \item Arrive home and drop the bag.
    \item Drink a glass of water and eat a snack.
    \item Study math for 45 minutes.
    \item Take a 15-minute break.
    \item Go to the grocery store for milk.
\end{enumerate}

This exact sequence of steps is called an **algorithm**.

\begin{definition}[Algorithm]
An algorithm is a finite sequence of clear, step-by-step instructions used to solve a problem or complete a task.
\end{definition}

In this setup, Alex acts as the ``processor'' following the instructions. But what if unexpected things happen—like discovering it is already 5:00 PM and the grocery store is closing? Alex must update the plan using a **decision**.

\subsection{Visualizing Decisions with a Flowchart}

When plans change, our logic changes too. We can draw this using a flowchart.

\begin{figure}[h]
    \centering
    \begin{tikzpicture}[node distance=2cm]
        \node (start) [startstop] {Start};
        \node (home) [io, below of=start] {Arrive Home};
        \node (snack) [process, below of=home] {Eat Snack};
        \node (decide) [decision, below of=snack, yshift=-0.8cm] {Is it 5:00 PM?};
        \node (study) [process, below of=decide, yshift=-0.8cm] {Study Math};
        \node (store) [process, right of=decide, xshift=3.5cm] {Go to Store};
        \node (end) [startstop, below of=study] {End};

        \draw [arrow] (start) -- (home);
        \draw [arrow] (home) -- (snack);
        \draw [arrow] (snack) -- (decide);
        \draw [arrow] (decide) -- node[anchor=east] {No} (study);
        \draw [arrow] (decide) -- node[anchor=south] {Yes} (store);
        \draw [arrow] (study) -- (end);
        \draw [arrow] (store) |- (end);
    \end{tikzpicture}
    \caption{Flowchart for Alex's Afternoon Decision Making}
    \label{fig:afternoon_flowchart}
\end{figure}

\subsection{Python Code Example: Printing a Daily Schedule}

Let us see how Alex can tell a computer to display this daily plan using Python.

\begin{lstlisting}
# Alex's Daily Planner Program
print("Step 1: Arrive home from school.")
print("Step 2: Eat a healthy snack.")
print("Step 3: Study math for 45 minutes.")
print("Step 4: Take a short break.")
print("Afternoon plan complete!")
\end{lstlisting}

\begin{pythonoutput}
Step 1: Arrive home from school.
Step 2: Eat a healthy snack.
Step 3: Study math for 45 minutes.
Step 4: Take a short break.
Afternoon plan complete!
\end{pythonoutput}

\textbf{Explanation of the Code:}
\begin{itemize}
    \item Lines starting with \texttt{\#} are **comments**. Python ignores them; they are just notes for human readers.
    \item \texttt{print(...)} is a built-in Python command that displays whatever text is inside the quotation marks onto the computer screen.
\end{itemize}

---

\section{Programming Without a Computer}

You do not need an electronic computer to practice programming. You use programming logic every day in your normal routines:
\begin{itemize}
    \item \textbf{Cooking a recipe:} Following ingredients and steps in order (An Algorithm).
    \item \textbf{Packing your school bag:} Checking tomorrow's timetable and packing books (Conditional logic).
    \item \textbf{Brushing your teeth:} Repeating a brushing motion back and forth twenty times (A Loop).
\end{itemize}

\section{Real-Life Problem: Managing Pocket Money}

Maya receives $500$ Taka as weekly pocket money. Maya wants to buy a book for $300$ Taka and save the rest. How does Maya calculate the remaining balance?

\begin{itemize}
    \item \textbf{Input:} Initial money = $500$, Book cost = $300$.
    \item \textbf{Process:} Subtract the book cost from the initial money ($500 - 300$).
    \item \textbf{Output:} Remaining balance = $200$ Taka.
\end{itemize}

Computers excel at performing calculations like this instantly. Let's write a short Python program for Maya's pocket money.

\begin{lstlisting}
# Maya's Pocket Money Calculator
total_money = 500
book_cost = 300

remaining_money = total_money - book_cost
print("Remaining Pocket Money:", remaining_money, "Taka")
\end{lstlisting}

\begin{pythonoutput}
Remaining Pocket Money: 200 Taka
\end{pythonoutput}

\textbf{Explanation of the Code:}
\begin{itemize}
    \item \texttt{total\_money = 500} stores the number $500$ inside a storage container named \texttt{total\_money}.
    \item \texttt{total\_money - book\_cost} subtracts the two numbers.
    \item \texttt{print(...)} outputs the final calculated result to the screen.
\end{itemize}

---

\section{The Input-Process-Output (IPO) Model}

Almost every computer program—and every daily task—follows the **IPO Model**. It consists of three simple stages:

\begin{table}[h]
    \centering
    \renewcommand{\arraystretch}{1.5}
    \begin{tabular}{|p{3cm}|p{6cm}|p{5cm}|}
        \hline
        \textbf{Component} & \textbf{Description} & \textbf{Real-Life Example} \\
        \hline
        \textbf{Input} & Raw data or information given to the system. & Typing \texttt{7:00 AM} into your phone alarm. \\
        \hline
        \textbf{Process} & The internal calculations or rules applied. & The phone's clock constantly checking the time. \\
        \hline
        \textbf{Output} & The final result shown to the user. & The phone ringing loudly at 7:00 AM. \\
        \hline
    \end{tabular}
    \caption{The Input-Process-Output (IPO) Model}
    \label{tab:ipo}
\end{table}

\subsection{Daily Example: Setting a Phone Alarm}

When Alex sets a morning alarm, the phone acts as a tiny computer following an IPO loop.

\begin{figure}[h]
    \centering
    \begin{tikzpicture}[node distance=1.5cm]
        \node (in) [io] {Input: Set alarm for 07:00 AM};
        \node (pro1) [process, below of=in] {Process: Check current time};
        \node (dec) [decision, below of=pro1, yshift=-0.8cm] {Is time 07:00 AM?};
        \node (wait) [process, left of=dec, xshift=-2.5cm] {Wait / Loop};
        \node (out) [io, below of=dec, yshift=-0.8cm] {Output: Alarm Rings!};

        \draw [arrow] (in) -- (pro1);
        \draw [arrow] (pro1) -- (dec);
        \draw [arrow] (dec) -- node[anchor=south] {Yes} (out);
        \draw [arrow] (dec) -| node[anchor=south] {No} (wait);
        \draw [arrow] (wait) -- (pro1);
    \end{tikzpicture}
    \caption{Phone Alarm Logic Flow}
    \label{fig:alarm_flow}
\end{figure}

\section{Why Computers Need Programmers}

Computers are lightning-fast, but they are also completely ``blind and deaf'' on their own. A computer cannot think, feel, or decide what to do next without a human programmer. It simply executes instructions mechanically.

If you give a computer bad instructions, it will make mistakes instantly. If you forget to tell a repeating task when to stop, the computer will loop forever (an infinite loop). Programming is how we bridge human creativity with machine speed!

---

\begin{exercise}
\begin{enumerate}
    \item Write down an everyday algorithm you follow every morning before leaving for school.
    \item Explain the IPO model using the example of withdrawing money from an ATM machine.
    \item What is the main difference between a human brain and a computer processor?
\end{enumerate}
\end{exercise}